\section{Improper integrals}
\label{sec:improper-integrals}
An improper integral is a \emph{limit of ordinary integrals}; it is not
an instruction to substitute an infinite number into an antiderivative.
For example,
\[
 \int_a^b f=\lim_{d\to b^-}\int_a^d f
\]
when the left endpoint is harmless. If both endpoints are troublesome,
choose an interior point $\alpha$ and take the two limits separately.
A finite limit means convergence. A limit $+\infty$ or $-\infty$ is an
extended value, not convergence. An oscillating limit has no value.
In particular, a symmetric limit on $[-R,R]$ cannot replace the two
one-sided limits in the definition.

For Stieltjes integrals, $F$ is increasing. Its increment
$F(v)-F(u)$ replaces the interval length $v-u$. We use Ross's
endpoint convention from Section~35, so that adjacent integrals add
without counting a jump twice. A distribution function has total
increment $F(\infty)-F(-\infty)=1$.

\exercise{36.1}{Agreement with an ordinary integral}
\goal If $f$ is Riemann integrable on the finite interval $[a,b]$,
prove that $\lim_{d\to b^-}\int_a^d f=\int_a^b f$.
\idea The only omitted contribution is over the short interval $[d,b]$.
\step{Bound that contribution.}
Riemann integrability includes boundedness. Choose $B\geq0$ with
$|f(x)|\leq B$ on $[a,b]$. Restrictions of an integrable function are
integrable, and additivity and the integral inequality give
\[
 \left|\int_a^b f-\int_a^d f\right|
 =\left|\int_d^b f\right|\leq\int_d^b |f|
 \leq B(b-d).
\]
If $B=0$, the difference is already zero. If $B>0$, given
$\varepsilon>0$, choosing $0<b-d<\varepsilon/B$ makes the difference
less than $\varepsilon$.
\answer{The difference tends to zero, so the two definitions agree.}

\exercise{36.2}{The splitting point does not change the value}
\goal Show that the definition of an integral on $(a,b)$ is independent
of the chosen interior point.
\step{Compare two splitting points by a finite integral.}
Suppose $a<\alpha<\beta<b$ and set
$C=\int_\alpha^\beta f$. This is a finite ordinary integral by the
local integrability hypothesis. For $a<c<\alpha$ and $\beta<d<b$,
\[
 \int_c^\beta f=\int_c^\alpha f+C,
 \qquad
 \int_\beta^d f=\int_\alpha^d f-C.
\]
Pass to the endpoint limits. Adding or subtracting a finite number
preserves the existence and sign of an infinite limit. Thus, if
$A=\int_a^\alpha f$ and $B=\int_\alpha^b f$ exist, then
\[
 \int_a^\beta f=A+C,\qquad \int_\beta^b f=B-C.
\]
\step{Check that the permitted sum is unchanged.}
If $A,B$ are finite, their sum is plainly unchanged. If one is infinite
and their sum is permitted, the finite changes do not alter either
infinite sign. An excluded sum $+\infty+(-\infty)$ remains excluded.
The same argument in reverse proves the converse, and interchanging
$\alpha,\beta$ covers either order.
\answer{The existence and the value of the improper integral do not
depend on the split.}

\exercise{36.3}{Power singularities at zero and at infinity}
\goal Analyze $\int_0^1 x^{-p}\,dx$ and $\int_0^\infty x^{-p}\,dx$
for $p>0$.
\step{(a) Cut off the singular endpoint before integrating.}
For $0<c<1$ and $p\ne1$, the power rule gives
\[
 \int_c^1 x^{-p}\,dx=\frac{1-c^{1-p}}{1-p}.
\]
If $0<p<1$, the exponent $1-p$ is positive, so $c^{1-p}\to0$;
the answer is $1/(1-p)$. If $p>1$, then
$c^{1-p}=1/c^{p-1}\to+\infty$. Rewriting the last display as
$(c^{1-p}-1)/(p-1)$ shows that the integral tends to $+\infty$.
For completeness, the boundary case $p=1$ has
$\int_c^1 dx/x=-\log c\to+\infty$.
\step{(b) Treat both ends separately.}
For $p\ne1$,
\[
 \int_1^d x^{-p}\,dx=\frac{d^{1-p}-1}{1-p}.
\]
This tends to $1/(p-1)$ if $p>1$, and to $+\infty$ if $0<p<1$.
For $p=1$ it equals $\log d\to+\infty$. Thus $p>1$ fails at
zero, $0<p<1$ fails at infinity, and $p=1$ fails at both ends.
All integrands are nonnegative, so there is no cancellation issue.
\answer{The integral on $(0,1)$ converges exactly for $0<p<1$;
the integral on $(0,\infty)$ is $+\infty$ for every $p>0$.}

\exercise{36.4}{Three explicit improper integrals}
\step{(a) The logarithm near zero.}
Integration by parts on $[c,1]$, where $c>0$, yields
\[
 \int_c^1\log x\,dx=[x\log x-x]_c^1=-1-c\log c+c.
\]
To justify $c\log c\to0$, put $t=1/c$. The elementary estimate
$\log t=\int_1^t du/u\leq\int_1^t du/\sqrt u
=2(\sqrt t-1)$ for $t\geq1$ gives
$0\leq(\log t)/t\leq2/\sqrt t\to0$.
Consequently $\int_0^1\log x\,dx=-1$.
\step{(b) A logarithm divided by its argument.}
For $d>2$, substitution $u=\log x$, $du=dx/x$, gives
\[
 \int_2^d\frac{\log x}{x}\,dx
 =\frac{(\log d)^2-(\log2)^2}{2}\longrightarrow+\infty.
\]
\step{(c) The reciprocal quadratic.}
Since $(\arctan x)'=1/(1+x^2)$,
\[
 \int_0^d\frac{dx}{1+x^2}=\arctan d\longrightarrow\frac\pi2.
\]
\answer{The values, in order, are $-1$, $+\infty$, and $\pi/2$.}

\exercise{36.5}{A nonnegative integral as a supremum}
\goal For continuous $f\geq0$ on $(a,b)$, identify its improper
integral with the supremum of all integrals over compact subintervals.
\step{Establish the two one-sided limits.}
Choose $\alpha\in(a,b)$. Define
\[
 A(c)=\int_c^\alpha f\quad(a<c<\alpha),\qquad
 B(d)=\int_\alpha^d f\quad(\alpha<d<b).
\]
Because $f\geq0$, $A(c)$ increases as $c$ moves toward $a$,
and $B(d)$ increases as $d$ moves toward $b$. The monotone-limit
property therefore gives limits $A,B\in[0,\infty]$, equal to their
respective suprema. The improper integral exists and equals $A+B$;
this sum can never have the forbidden form $+\infty+(-\infty)$.
\step{Compare with all compact subintervals.}
Let $M=\sup\{\int_c^d f:[c,d]\subset(a,b)\}$. Enlarging any
$[c,d]$ so that it contains $\alpha$ can only increase its integral.
An enlarged interval has integral $A(c')+B(d')\leq A+B$; thus
$M\leq A+B$. Conversely, all intervals with $c<\alpha<d$ are
allowed in the supremum, so $M\geq A(c)+B(d)$. If $A,B$ are
finite, choose $c,d$ making these terms within $\varepsilon/2$ of
their suprema, and obtain $M\geq A+B-\varepsilon$. If either
supremum is infinite, these same integrals exceed every prescribed
bound, so $M=+\infty$.
\answer{$\displaystyle\int_a^b f=\sup_{[c,d]\subset(a,b)}\int_c^d f$,
with $+\infty$ allowed.}

\exercise{36.6}{Comparison tests}
\goal Suppose $0\leq f\leq g$ and both functions are continuous on
$(a,b)$. Prove the convergence and divergence comparisons.
\step{Compare finite intervals first.}
For every compact $[c,d]\subset(a,b)$, ordinary integral monotonicity
gives $0\leq\int_c^d f\leq\int_c^d g$. Exercise~36.5 allows us to
take suprema, giving
\[
 0\leq\int_a^b f\leq\int_a^b g\leq+\infty.
\]
\step{(a) Finite upper bound.}
If $\int_a^b g$ is finite, the middle quantity is bounded above by
that finite number. The nonnegative integral of $f$ therefore converges.
\step{(b) Infinite lower bound.}
If $\int_a^b f=+\infty$, no finite number can bound the integral
of $g$ from above. Therefore $\int_a^b g=+\infty$.
\answer{A smaller nonnegative function inherits convergence; a larger
one inherits divergence to $+\infty$.}

\exercise{36.7}{The Gaussian integral}
\step{(a) Prove finiteness before attempting evaluation.}
The integrand is even. On $[0,1]$ it is at most $1$; on $[1,\infty)$,
$x^2\geq x$, so $e^{-x^2}\leq e^{-x}$. The comparison test gives
\[
 0<I:=\int_{-\infty}^{\infty}e^{-x^2}\,dx
 \leq2\left(1+\int_1^\infty e^{-x}\,dx\right)
 =2(1+e^{-1})<\infty.
\]
\step{(b) Explain the double-integral calculation.}
This part uses polar coordinates for a double integral, as requested
by the exercise. The following bounded-region argument avoids assuming
in advance that improper iterated integrals may be interchanged.
Set $I_R=\int_{-R}^R e^{-x^2}\,dx$. On the square
$Q_R=[-R,R]^2$, integration of the continuous product gives
\[
 \iint_{Q_R}e^{-x^2-y^2}\,dx\,dy=I_R^2.
\]
Let $D_R$ be the disk of radius $R$. The inclusions
$D_R\subset Q_R\subset D_{\sqrt2R}$ and positivity give a lower
and an upper bound for $I_R^2$. In polar coordinates,
\[
 \iint_{D_R}e^{-x^2-y^2}\,dx\,dy
 =\int_0^{2\pi}\int_0^R e^{-r^2}r\,dr\,d\theta
 =\pi(1-e^{-R^2}).
\]
Here the extra factor $r$ is the polar area factor, and
$\int_0^R e^{-r^2}r\,dr=(1-e^{-R^2})/2$. Hence
\[
 \pi(1-e^{-R^2})\leq I_R^2\leq\pi(1-e^{-2R^2}).
\]
Both outer expressions tend to $\pi$. Since part (a) shows
$I_R\to I$, we obtain $I^2=\pi$; and $I>0$ selects the positive root.
\answer{$\displaystyle\int_{-\infty}^{\infty}e^{-x^2}\,dx=\sqrt\pi$.}

\exercise{36.8}{Absolute convergence implies convergence}
\goal If $f$ is continuous on $(a,b)$ and $\int_a^b|f|<\infty$,
prove that $\int_a^b f$ exists and is finite.
\idea Write $f$ as the difference of two nonnegative functions whose
integrals are bounded above by integrals already known to converge.
\step{Use positive and negative parts.}
Define
\[
 f_+(x)=\max\{f(x),0\},\qquad
 f_-(x)=\max\{-f(x),0\}.
\]
These are continuous because, for example,
$f_+=(|f|+f)/2$. Pointwise, $0\leq f_+,f_-\leq|f|$ and
$f=f_+-f_-$. Exercise~36.6 proves that both nonnegative improper
integrals are finite.
\step{Pass to limits without subtracting infinities.}
On each compact subinterval, ordinary linearity gives
$\int_c^d f=\int_c^d f_+-\int_c^d f_-$. Fix an interior split,
and take the left and right endpoint limits separately. Each term
has a finite limit, so the limit law for subtraction applies.
\answer{$\int_a^b f$ converges and
$\left|\int_a^b f\right|\leq\int_a^b|f|$.}

\exercise{36.9}{Integrals for the normal distribution}
Write $\varphi(x)=(2\pi)^{-1/2}e^{-x^2/2}$ and
$F(t)=\int_{-\infty}^t\varphi(x)\,dx$. Substitution
$x=\sqrt2u$ in Exercise~36.7 gives $\int_\R\varphi=1$.
\step{(a) Replace the Stieltjes integrator by its density.}
The Fundamental Theorem of Calculus gives $F'=\varphi$, a continuous
function. For continuous $f$, Theorem~35.13 on each finite $[a,b]$
therefore yields
\[
 \int_a^b f\,dF=\int_a^b f(x)F'(x)\,dx
 =\frac1{\sqrt{2\pi}}\int_a^b f(x)e^{-x^2/2}\,dx.
\]
Taking the two endpoint limits separately proves the same identity
on $\R$ whenever the integral specified in the question exists;
positive constant multiplication preserves finite and infinite limits.
\step{(b) The second moment.}
Since $\varphi'(x)=-x\varphi(x)$, integration by parts on $[a,b]$ gives
\[
 \int_a^b x^2\varphi(x)\,dx
 =[-x\varphi(x)]_a^b+\int_a^b\varphi(x)\,dx.
\]
The boundary term tends to zero at either infinity. For instance,
$e^t\geq t^2/2$ for $t\geq0$ implies
$xe^{-x^2/2}\leq8/x^3$ for $x>0$, and negative $x$ are handled
by absolute values. Thus $\int_\R x^2\,dF=1$.
\step{(c) The exponential of the square.}
The integrand here is $e^{x^2}$, not $e^x$. Its density product is
\[
 e^{x^2}\varphi(x)=\frac1{\sqrt{2\pi}}e^{x^2/2}
 \geq\frac1{\sqrt{2\pi}}.
\]
The integral of this positive constant over an infinite half-line
is infinite, so comparison proves $\int_\R e^{x^2}\,dF=+\infty$.
\step{(d) The absolute first moment.}
Evenness and the antiderivative of $x e^{-x^2/2}$ give
\[
 \int_\R |x|\,dF
 =\frac2{\sqrt{2\pi}}\int_0^\infty x e^{-x^2/2}\,dx
 =\frac2{\sqrt{2\pi}}[-e^{-x^2/2}]_0^\infty
 =\sqrt{\frac2\pi}.
\]
\step{(e) The first moment and a necessary convergence check.}
Part (d) shows absolute convergence. Each half-integral is therefore
finite. Because $x\varphi(x)$ is odd, substitution $u=-x$ in the
negative half-integral makes it the negative of the positive half.
Their sum is zero. Oddness alone would not suffice without the
existence of both halves.
\answer{The values in (b)--(e) are $1$, $+\infty$,
$\sqrt{2/\pi}$, and $0$.}

\exercise{36.10}{Linearity on the whole real line}
\goal For $F$-integrable $f,g$ on $\R$, prove integrability and
the sum rule for $f+g$.
\step{Apply finite-interval linearity.}
For every finite $a<0<b$, Theorem~35.8 gives
\[
 \int_a^0(f+g)\,dF=\int_a^0 f\,dF+\int_a^0 g\,dF,
 \qquad
 \int_0^b(f+g)\,dF=\int_0^b f\,dF+\int_0^b g\,dF.
\]
Since $f,g$ are $F$-integrable on $\R$, all four half-line limits
are finite. The limit law for addition proves that the two limits
for $f+g$ exist and are finite. Add them and regroup finite real numbers.
\answer{$f+g$ is $F$-integrable and
$\int_\R(f+g)\,dF=\int_\R f\,dF+\int_\R g\,dF$. The same
argument proves scalar multiplication, including multiplication by $-1$.}

\exercise{36.11}{Monotonicity on the whole real line}
\goal If $f,g$ are $F$-integrable and $f\leq g$, prove the integral
inequality.
\step{Preserve the inequality through finite limits.}
Theorem~35.9 gives
$\int_a^0 f\,dF\leq\int_a^0 g\,dF$ for each $a<0$, and
$\int_0^b f\,dF\leq\int_0^b g\,dF$ for each $b>0$.
Both sides have finite limits by hypothesis. If real-valued
quantities satisfy an inequality before taking a limit, the limits
satisfy the same weak inequality. Pass to $a\to-\infty$ and
$b\to+\infty$ separately, and add.
\answer{$\displaystyle\int_\R f\,dF\leq\int_\R g\,dF$.}

\exercise{36.12}{Comparison for improper Stieltjes integrals}
\goal State and prove the analogue of Exercise~36.6.
Assume $F$ is increasing and $f,g$ are $F$-integrable on every
finite interval, with $0\leq f\leq g$. Continuity of $f,g$ is
one sufficient condition for local integrability.
\step{Nonnegativity supplies the endpoint limits.}
For $b>0$, $\int_0^b f\,dF$ is increasing in $b$, since its
increment over an added interval is nonnegative. Likewise,
$\int_a^0 f\,dF$ increases as $a\to-\infty$. Thus both limits
exist in $[0,\infty]$. The same holds for $g$.
\step{Compare before and after the limits.}
Finite-interval monotonicity gives
\[
 0\leq\int_a^0 f\,dF\leq\int_a^0 g\,dF,
 \qquad
 0\leq\int_0^b f\,dF\leq\int_0^b g\,dF.
\]
Taking limits and adding nonnegative extended values proves
$0\leq\int_\R f\,dF\leq\int_\R g\,dF$.
\answer{If $\int_\R g\,dF<\infty$, then $\int_\R f\,dF<\infty$.
If $\int_\R f\,dF=+\infty$, then $\int_\R g\,dF=+\infty$.}

\exercise{36.13}{Absolute convergence for Stieltjes integrals}
\goal If $f$ is continuous on $\R$ and $\int_\R|f|\,dF<\infty$,
prove $f$ is $F$-integrable on $\R$.
\step{Obtain finite nonnegative integrals.}
Continuity makes $f$, $|f|$, and $h=f+|f|$ locally $F$-integrable.
Pointwise, $0\leq h\leq2|f|$. Exercise~36.12 shows that
$h$ has a finite integral, since $2|f|$ does.
\step{Subtract a function with a known finite integral.}
Scalar multiplication shows $-|f|$ is $F$-integrable.
Exercise~36.10 now applies to $f=h+(-|f|)$, so $f$ is
$F$-integrable. Applying monotonicity to $-|f|\leq f\leq|f|$
also gives $|\int_\R f\,dF|\leq\int_\R|f|\,dF$.
\answer{Absolute $F$-integrability implies $F$-integrability. The
argument still works if continuity is replaced by local
$F$-integrability of $f$ and $|f|$.}

\exercise{36.14}{A countable collection of point masses}
\goal Let the distinct points be $u_j\in\R$, let $c_j\geq0$, and
suppose $C=\sum_{j=1}^\infty c_j<\infty$. Analyze
$F(t)=\sum_j c_j J_{u_j}(t)$, where $J_u(t)=0$ for $t<u$ and
$J_u(t)=1$ for $t\geq u$ on $\R$.
\step{(a) Existence and monotonicity.}
For each $t$, $0\leq c_jJ_{u_j}(t)\leq c_j$. The series therefore
converges, and its tail is at most $\sum_{j>N}c_j$, independently
of $t$. Thus convergence is uniform. Each summand is increasing;
for $s<t$, termwise comparison gives $F(s)\leq F(t)$.
\step{(b) Finite intervals and their endpoint convention.}
Let $|f|\leq B$ on $\R$. Theorem~35.12(b), or its finite-mass
tail estimate, makes every bounded $f$ $F$-integrable on $[a,b]$.
For this right-continuous $F$, the formula is
\[
 \int_a^b f\,dF=\sum_{j:\,a<u_j\leq b}c_j f(u_j).
\]
Indeed a single global step contributes its mass exactly when
$a<u_j\leq b$; a jump at $a$ has already occurred before the
interval's increment begins. Summing the steps is justified because
the absolute value of the contribution from all indices $j>N$ is
at most $B\sum_{j>N}c_j$. This tends to zero.
\step{Pass to the whole line, controlling the tails.}
The numerical series $\sum_j c_jf(u_j)$ is absolutely convergent,
because $\sum_j|c_jf(u_j)|\leq BC$. To see directly that the
finite-interval sums approach it, choose $N$ with
$B\sum_{j>N}c_j<\varepsilon$ (the case $B=0$ is immediate),
then choose $a<\min_{j\leq N}u_j$ and
$b>\max_{j\leq N}u_j$. All first $N$ terms are included, and
the omitted terms have total absolute value below $\varepsilon$.
The same argument on the two halves treats $u_j=0$ in the negative
half only. Hence
\[
 \int_{-\infty}^{\infty}f\,dF=\sum_{j=1}^\infty c_jf(u_j).
\]
\step{(c) The total mass condition.}
The same finite-head/small-tail argument gives
$F(-\infty)=0$ and $F(\infty)=C$. Consequently $F$ is a
distribution function exactly when $\sum_jc_j=1$. It is also
right-continuous: each finite partial sum is right-continuous,
and uniform convergence preserves this one-sided continuity.
\step{(d) Dense jumps and a necessary qualification.}
For $s<t$,
\[
 F(t)-F(s)=\sum_{j:\,s<u_j\leq t}c_j.
\]
If the $u_j$ enumerate $\Q$ and \emph{every $c_j>0$}, there is a
rational $u_j$ strictly between $s$ and $t$, so this difference is
positive. Thus $F$ is strictly increasing. More generally the exact
condition is that $\{u_j:c_j>0\}$ is dense in $\R$.
\answer{Parts (a)--(c) hold for nonnegative coefficients as stated.
Part (d), read with merely nonnegative coefficients, needs the
additional density condition: all $c_j=0$ would give $F=0$.
This is the positive-coefficient situation of Exercise~35.12.}

\exercise{36.15}{Uniform convergence on spaces of infinite and finite mass}
\step{(a) Small height can be offset by large width.}
Define $f_n(x)=1/n$ on $[0,n]$, and $f_n(x)=0$ elsewhere.
Each function is piecewise constant and has a finite improper
Riemann integral. Its rectangle has height $1/n$ and width $n$, so
\[
 \int_\R f_n(x)\,dx=1,\qquad
 \sup_{x\in\R}|f_n(x)-0|=\frac1n\longrightarrow0.
\]
Thus uniform convergence alone does not preserve an integral over
a domain of infinite ordinary length.
\step{(b) Local integrability of the uniform limit.}
Now suppose $F$ is a distribution function, all $f_n$ are
$F$-integrable on $\R$, and $f_n\to f$ uniformly on $\R$.
On each finite $[a,b]$, the uniform-limit theorem for Stieltjes
integrals (Exercise~35.6) makes $f$ $F$-integrable there.
For clarity, that theorem follows from
\[
 U_F(f,P)-L_F(f,P)
 \leq U_F(f_n,P)-L_F(f_n,P)
       +2\|f-f_n\|_\infty [F(b)-F(a)]:
\]
first choose $n$ to make the second term small, then choose $P$
using integrability of $f_n$.
\step{Prove global integrability without assuming global boundedness.}
Choose one index $N$ with $\|f-f_N\|_\infty\leq1$. The difference
$h=f-f_N$ is bounded and locally $F$-integrable. Since $F$ has
finite total increment, Theorem~36.4 makes $h$ $F$-integrable
on $\R$. Exercise~36.10 now shows that $f=f_N+h$ is
$F$-integrable. This argument does not require $f_N$, or $f$,
to be bounded on all of $\R$.
\step{Estimate the difference of the integrals.}
For any sufficiently large $n$, put
$\eta_n=\|f-f_n\|_\infty$. By monotonicity and the total-mass
identity $\int_\R1\,dF=F(\infty)-F(-\infty)=1$,
\[
 -\eta_n\leq\int_\R(f-f_n)\,dF\leq\eta_n.
\]
Linearity therefore gives
\[
 \left|\int_\R f\,dF-\int_\R f_n\,dF\right|
 \leq\eta_n\longrightarrow0.
\]
\answer{For a distribution function, uniform convergence does preserve
the integral. The decisive bound is the finite total mass $1$,
rather than the geometric length of the real line.}
